% Illustrative Math 1151-style example using the current ximeraExam commands.
% This is a NEW example, not a transcription of an existing course exam.
% Keep ximeraExam.sty beside this file and compile with LuaLaTeX twice.
\DocumentMetadata{lang=en-US,pdfstandard=ua-2,tagging=on}
\documentclass{ximera}
\usepackage{unicode-math}
\usepackage[points]{ximeraExam}

\examtitle{Midterm 1 Worked Example}
\examcourse{MATH 101}
\examdate{Sample Date}
\examversion{Form A}


% Customize the cover page's header and footer, if you enable its page style.
% Each command takes left, center, and right fields, in that order.
\firstpageheader{}{\examtitletext}{}
\firstpagefooter{}{\examcoursetext}{}

% Customize the header and footer on the exam's question pages.
% Leave a field empty ({}) to print nothing in that position.
\runningheader{\examcoursetext}{\examtitletext}{\examversiontext}
\runningfooter{\examdatetext}{}{Page~\thepage\ of~\pageref{LastPage}}

% Switch to \printanswers to see solutions and remove student response space.
\printanswers

\begin{document}
\begin{examcover}
  \examcoverheading{\examtitletext}
  \noindent\examcoursetext\quad\examversiontext\quad\examdatetext\par
  \noindent\examnameprompt\ \examnamefield\par\bigskip
  \noindent\textbf{Instructions:} Give exact answers and show supporting work.
  \par\bigskip
  \gradetable[h][][4]
\end{examcover}

\begin{questions}

% A two-part limit question demonstrates point inheritance.
\question Evaluate the following limits, showing your reasoning.
\begin{parts}
  \part[4] \(\displaystyle\lim_{x\to 2}\frac{x^2-4}{x-2}\).
  \begin{solution}[1.25in]
    Factor and cancel for \(x\ne2\):
    \(\frac{(x-2)(x+2)}{x-2}=x+2\). The limit is \(4\).
  \end{solution}

  \part[4] \(\displaystyle\lim_{x\to 0}\frac{\sin x}{x}\).
  \begin{solution}[1.25in]
    The standard trigonometric limit is \(1\).
  \end{solution}
\end{parts}

% The same spacing syntax works for a longer written justification.
\question[6] Let \(f(x)=x^2-2\). Explain why \(f\) has a zero in \((1,2)\).
\begin{solution}[\stretch{1}]
  Polynomials are continuous. Since \(f(1)=-1\) and \(f(2)=2\),
  the Intermediate Value Theorem guarantees a number \(c\in(1,2)\)
  for which \(f(c)=0\).
\end{solution}

% A short final-answer line can be used alongside a worked solution.
\question[4] Differentiate \(h(x)=x^3-4x\).
\begin{solution}[1in]
  Using the power rule, \(h'(x)=3x^2-4\).
\end{solution}

\end{questions}
\end{document}
